% GENERATED FILE. Edit canonical lessons, then run npm run generate:books.
% canonical-source-hash: fnv1a64:fda5d1d18642b252
% canonical-lessons: FR-C287-cambrioler, FR-C287-cambrioleur, FR-C287-cambriolage, FR-C287-effraction, FR-C287-butin

\chapter{Le cambriolage}
\label{ch:fr-le-cambriolage}
\hlchaptermodality{287}

\begin{chapteropening}
\textbf{By the end of this chapter:} \emph{I can tell the story of a burglary: breaking in, the burglar, the burglary, the forced entry and the loot.}

Complete the last lesson of chapter 287: I can tell the story of a burglary: breaking in, the burglar, the burglary, the forced entry and the loot.
\end{chapteropening}

\section[cambrioler]{cambrioler — to burgle, to break into}

\label{lesson:FR-C287-cambrioler}



\begin{quote}
Before the new one: say the French for the drop, then the French for the statistic.
\end{quote}

\subsection*{You'll want to know: cambrioler}
\textbf{cambrioler} — \textquotedblleft{}to burgle, to break into\textquotedblright{}.

\begin{quote}
Quelqu'un a cambriolé la maison pendant les vacances.
\end{quote}

\begin{quote}
\emph{Someone burgled the house during the holidays.}
\end{quote}

Said of a place, not a person: \emph{cambrioler un appartement}.

\begin{cousinweb}
\emph{cambrioler} is built on the old slang word \emph{cambriole} \textquotedblleft{}a small room\textquotedblright{}, usually traced to Provençal \emph{cambro} \textquotedblleft{}room\textquotedblright{}, which goes back to Latin \emph{camera} \textquotedblleft{}a vault\textquotedblright{} (Lewis \& Short); EtymDB lists the verb without a parent.
\end{cousinweb}

\begin{grammarlens}[title={What you've built}]
One more word for telling, explaining and describing at B1.
\end{grammarlens}

\subsection*{Your turn}
\begin{itemize}
\raggedright
  \item \emph{Say it:} \emph{cambrioler}
  \item \emph{Say it:} \emph{cambrioler}, once more
  \item \emph{Say it:} a sentence of your own with \emph{cambrioler}
  \item \emph{Recall:} say the French for the drop, then the French for the statistic, then say \emph{cambrioler} again
\end{itemize}

\begin{tcolorbox}[breakable,colback=teal!4,colframe=teal!35!black,title={Before you move on}]
What does cambrioler mean? (\textquotedblleft{}To burgle, to break into\textquotedblright{}.) Say it once more.
\end{tcolorbox}
\section[le cambrioleur]{le cambrioleur — the burglar}

\label{lesson:FR-C287-cambrioleur}



\begin{quote}
Before the new one: say the French for the statistic, then the French for to burgle.
\end{quote}

\subsection*{You'll want to know: le cambrioleur}
\textbf{le cambrioleur} — \textquotedblleft{}the burglar\textquotedblright{}.

\begin{quote}
Le cambrioleur est entré par la fenêtre.
\end{quote}

\begin{quote}
\emph{The burglar got in through the window.}
\end{quote}

Feminine: \emph{la cambrioleuse}.

\begin{cousinweb}
\emph{cambrioleur} is built on \emph{cambrioler} \textquotedblleft{}to burgle\textquotedblright{} with the agent suffix \emph{-eur} (EtymDB); the verb goes back to slang \emph{cambriole} \textquotedblleft{}a small room\textquotedblright{}.
\end{cousinweb}

\begin{grammarlens}[title={What you've built}]
One more word for telling, explaining and describing at B1.
\end{grammarlens}

\subsection*{Your turn}
\begin{itemize}
\raggedright
  \item \emph{Say it:} \emph{le cambrioleur}
  \item \emph{Say it:} \emph{le cambrioleur}, once more
  \item \emph{Say it:} a sentence of your own with \emph{le cambrioleur}
  \item \emph{Recall:} say the French for the statistic, then the French for to burgle, then say \emph{le cambrioleur} again
\end{itemize}

\begin{tcolorbox}[breakable,colback=teal!4,colframe=teal!35!black,title={Before you move on}]
What does le cambrioleur mean? (\textquotedblleft{}The burglar\textquotedblright{}.) Say it once more.
\end{tcolorbox}
\section[le cambriolage]{le cambriolage — the burglary, the break-in}

\label{lesson:FR-C287-cambriolage}



\begin{quote}
Before the new one: say the French for to burgle, then the French for the burglar.
\end{quote}

\subsection*{You'll want to know: le cambriolage}
\textbf{le cambriolage} — \textquotedblleft{}the burglary, the break-in\textquotedblright{}.

\begin{quote}
Le cambriolage a eu lieu pendant la nuit.
\end{quote}

\begin{quote}
\emph{The burglary took place during the night.}
\end{quote}

Said of a home, a shop or an office.

\begin{cousinweb}
\emph{cambriolage} is built on \emph{cambrioler} \textquotedblleft{}to burgle\textquotedblright{} with the noun suffix \emph{-age} (EtymDB).
\end{cousinweb}

\begin{grammarlens}[title={What you've built}]
One more word for telling, explaining and describing at B1.
\end{grammarlens}

\subsection*{Your turn}
\begin{itemize}
\raggedright
  \item \emph{Say it:} \emph{le cambriolage}
  \item \emph{Say it:} \emph{le cambriolage}, once more
  \item \emph{Say it:} a sentence of your own with \emph{le cambriolage}
  \item \emph{Recall:} say the French for to burgle, then the French for the burglar, then say \emph{le cambriolage} again
\end{itemize}

\begin{tcolorbox}[breakable,colback=teal!4,colframe=teal!35!black,title={Before you move on}]
What does le cambriolage mean? (\textquotedblleft{}The burglary, the break-in\textquotedblright{}.) Say it once more.
\end{tcolorbox}
\section[l'effraction]{l'effraction — the forced entry, breaking in}

\label{lesson:FR-C287-effraction}



\begin{quote}
Before the new one: say the French for the burglar, then the French for the burglary.
\end{quote}

\subsection*{You'll want to know: l'effraction}
\textbf{l'effraction} — \textquotedblleft{}the forced entry, breaking in\textquotedblright{}.

\begin{quote}
La police n'a trouvé aucune trace d'effraction.
\end{quote}

\begin{quote}
\emph{The police found no sign of a forced entry.}
\end{quote}

\emph{Entrer par effraction} — to break in.

\begin{cousinweb}
\emph{effraction} is a learned word built on Latin \emph{effractus}, the past participle of \emph{effringere} \textquotedblleft{}to break open\textquotedblright{}, from \emph{ex} and \emph{frangere} \textquotedblleft{}to break\textquotedblright{} (Lewis \& Short).
\end{cousinweb}

\begin{grammarlens}[title={What you've built}]
One more word for telling, explaining and describing at B1.
\end{grammarlens}

\subsection*{Your turn}
\begin{itemize}
\raggedright
  \item \emph{Say it:} \emph{l'effraction}
  \item \emph{Say it:} \emph{l'effraction}, once more
  \item \emph{Say it:} a sentence of your own with \emph{l'effraction}
  \item \emph{Recall:} say the French for the burglar, then the French for the burglary, then say \emph{l'effraction} again
\end{itemize}

\begin{tcolorbox}[breakable,colback=teal!4,colframe=teal!35!black,title={Before you move on}]
What does l'effraction mean? (\textquotedblleft{}The forced entry, breaking in\textquotedblright{}.) Say it once more.
\end{tcolorbox}
\section[le butin]{le butin — the loot, the haul}

\label{lesson:FR-C287-butin}



\begin{quote}
Before the new one: say the French for the burglary, then the French for the forced entry.
\end{quote}

\subsection*{You'll want to know: le butin}
\textbf{le butin} — \textquotedblleft{}the loot, the haul\textquotedblright{}.

\begin{quote}
Le butin était caché dans une valise.
\end{quote}

\begin{quote}
\emph{The loot was hidden in a suitcase.}
\end{quote}

Also the haul of a happy day: \emph{le butin des soldes}.

\begin{cousinweb}
\emph{butin} comes from Old French \emph{butin} \textquotedblleft{}booty\textquotedblright{} (EtymDB), which is usually traced to Middle Low German \emph{būte} \textquotedblleft{}exchange, a share of booty\textquotedblright{}.
\end{cousinweb}

\begin{grammarlens}[title={What you've built}]
One more word for telling, explaining and describing at B1.
\end{grammarlens}

\subsection*{Your turn}
\begin{itemize}
\raggedright
  \item \emph{Say it:} \emph{le butin}
  \item \emph{Say it:} \emph{le butin}, once more
  \item \emph{Say it:} a sentence of your own with \emph{le butin}
  \item \emph{Recall:} say the French for the burglary, then the French for the forced entry, then say \emph{le butin} again
\end{itemize}

\begin{tcolorbox}[breakable,colback=teal!4,colframe=teal!35!black,title={Before you move on}]
What does le butin mean? (\textquotedblleft{}The loot, the haul\textquotedblright{}.) Say it once more.
\end{tcolorbox}
